\section{Master Structural Theorem of the Companion}
\label{sec:master-structural-theorem}


\begin{theorem}[Master structural theorem]
\label{thm:master-structural}
The Companion is governed by one common structural architecture with the
following ordered logic:
\[
\begin{aligned}
&\text{triolic source}
 \;\longrightarrow\;
 \text{admissible carrier}
 \;\longrightarrow\;
 \mathcal M_{16}
 \;\longrightarrow\;
 \mathcal H_{16}
 \;\longrightarrow\;
 \text{closure-stable sector}\\[0.2em]
&\longrightarrow\;\text{group-layer symmetry readout}
 \;\longrightarrow\;\text{null axis}\\[0.2em]
&\longrightarrow\;\text{photon fix-sector}
 \;\longrightarrow\;\text{residual localization}\\[0.2em]
&\longrightarrow\;\text{matter / tachyon / photon sector readout}\\[0.2em]
&\longrightarrow\;\text{Higgs mediation}
 \;\longrightarrow\;\text{field-equation readout}.
\end{aligned}
\]
\end{theorem}

\begin{proof}
The triolic source supplies the primitive structural object. The carrier layer
provides its admissible support, while the Millennium matrix transports this
same object through admissible realization changes. Under the
Heisenberg-compensator reading, the same matrix becomes the global balancing
mechanism that stabilizes the closure-compatible sector. The group layer then
supplies admissible symmetry actions on already given stabilized structure.
The null axis is the structurally distinguished sector carried through the
cascade, and its terminal compensated readout is the photon fix-sector.
Relative to that fix-sector, closure-stable residual localization yields the
matter branch, while its $J$-dual realization yields the tachyon branch. The
Higgs sector mediates between the compensated photon reference and the
persistence of residual localization. Finally, the effective field equations
read out the dynamical content of this already stabilized architecture.
\end{proof}

\begin{remark}[Compact roadmap for the reader]
The reader may keep the following compact roadmap in mind:
\[
\boxed{\begin{aligned}
&\text{one source} \to \text{one transported architecture}\\
&\to \text{one compensated closure logic}\\
&\to \text{three principal sector readouts} \to \text{one mediated field dynamics}
\end{aligned}}
\]
\end{remark}

\subsection{Schematic Master Diagram}
\label{subsec:master-diagram}

\begin{definition}[Schematic master diagram]
\label{def:master-diagram}
The structural logic of the Companion may be compressed into the following
schematic master diagram:
\[
\mathfrak T
\;\longrightarrow\;
\mathfrak C
\;\xrightarrow{\;\mathcal M_{16}\;}
\mathfrak C^{\mathrm{adm}}
\;\xrightarrow{\;\mathcal H_{16}\;}
\mathfrak C^{\mathrm{cl}}
\;\longrightarrow\;
\mathfrak G
\;\longrightarrow\;
\mathfrak L
\;\longrightarrow\;
\mathfrak P
\;\longrightarrow\;
\mathfrak R
\;\longrightarrow\;
\mathfrak H
\;\longrightarrow\;
\mathfrak F.
\]
Here:
\begin{align*}
\mathfrak T &:= \text{triolic source},\\
\mathfrak C &:= \text{carrier layer},\\
\mathfrak C^{\mathrm{adm}} &:= \text{admissibly transported carrier},\\
\mathfrak C^{\mathrm{cl}} &:= \text{closure-stable compensated carrier},\\
\mathfrak G &:= \text{group-layer symmetry readout},\\
\mathfrak L &:= \text{null-axis structure},\\
\mathfrak P &:= \text{photon fix-sector},\\
\mathfrak R &:= \text{residual localization sector},\\
\mathfrak H &:= \text{Higgs-mediation sector},\\
\mathfrak F &:= \text{field-equation readout}.
\end{align*}
\end{definition}

\begin{remark}[Language regime along the master chain]
\label{rem:language-regime-master-chain}
The Companion follows a structure-first logic along the master chain up to the
point at which a stable physical readout becomes available. Accordingly, the
earlier stages
\[
\mathfrak T
\to
\mathfrak C
\to
\mathcal M_{16}
\to
\mathcal H_{16}
\to
\mathfrak G
\to
\mathfrak L
\]
are described primarily in structurally prior language. By contrast, once the
readout has stabilized at the level of
\[
\mathfrak P
\to
\mathfrak R
\to
\mathfrak H
\to
\mathfrak F,
\]
physically explicit language is not only allowed but required. Thus the shift
from structural terminology to physical terminology is intentional and marks a
readout threshold, not a change of underlying object.
\end{remark}

\subsection{Matter, Tachyon, and Photon as Three Readouts of One Branch Architecture}
\label{subsec:mtp-readout}

\begin{definition}[Matter branch]
The matter branch is the closure-stable bound branch whose physical readout
retains nontrivial localization relative to the photon fix-sector.
\end{definition}

\begin{definition}[Tachyon branch]
The tachyon branch is the $J$-dual readout of the matter branch:
\[
K_2 = J\!\cdot\!K_1.
\]
It is the branch in which the same structural source is read in its frame-dual,
non-bound, or overextended transport mode relative to the compensated photon
sector.
\end{definition}

\begin{definition}[Photon branch]
The photon branch is the canonical compensated fix-sector readout $K_3$, i.e. the terminal null-axis sector selected by the Heisenberg compensator.
\end{definition}

\begin{theorem}[Three-readout principle]
\label{thm:three-readout-principle}
Matter, tachyon, and photon are three structurally linked readouts of one
common branch architecture: matter is the bound closure-stable localization
readout, tachyon is the $J$-dual unbound readout of the same branch source,
and photon is the compensated fix-sector readout.
\end{theorem}

\begin{proof}
Admissible sector transport does not create new ontology. The compensated
sector is the canonical fix-sector. Hence the matter branch is determined by
persistent closure-stable localization relative to this sector. The involution
$J$ generates the frame-dual branch of the same source, which yields the
tachyonic readout. The compensator selects the stabilized null/fix-sector,
which yields the photon readout. Therefore all three branches belong to one
common structural architecture and differ only by their mode of realization and
readout.
\end{proof}

\subsection{Mass Readout Relative to the Null Axis and Photon Fix-Sector}
\label{subsec:mass-null-axis}

\begin{definition}[Photon fix-sector revisited]
The photon sector is the canonical compensated readout of the global null
axis:
\[
L=\bigcap_{n\ge 0} L_n = S_0.
\]
It is the sector in which the frame-sensitive imbalance has been fully
neutralized and the structural readout lies entirely in the HC-fix sector.
\end{definition}

\begin{definition}[Massive branch]
A branch is called massive if it remains closure-stably localized and does not
collapse completely into the global null/fix sector. Equivalently, a massive
branch is one whose physical readout retains a nontrivial localization residue
relative to the photon sector.
\end{definition}

\begin{theorem}[Massive versus massless readout]
\label{thm:massive-vs-massless}
Within the MM/HC architecture, the distinction between massive and massless
physical readout is the distinction between complete reduction to the
compensated photon fix-sector and closure-stable persistence of localization
relative to that fix-sector.
\end{theorem}

\begin{proof}
If a branch reduces completely to the null/fix sector, no residual
localization remains, and the physical readout is massless. If, however, a
branch remains closure-stably localized and does not reduce completely to the
null/fix sector, then a nonzero localization residue survives relative to the
photon reference sector. This is precisely the massive readout.
\end{proof}

\subsection{Higgs Readout as Mediation Between Fix-Sector and Residual Localization}
\label{subsec:higgs-fix-localization}

\begin{definition}[Higgs mediation principle]
\label{def:higgs-mediation}
The Higgs readout is the effective mediation principle that governs how
closure-stable localization is maintained, transferred, or expressed relative
to the compensated photon fix-sector.
\end{definition}

\begin{theorem}[Structural role of the Higgs sector]
\label{thm:higgs-role-refined}
Within the MM/HC architecture, the Higgs sector is the distinguished coupling
readout that mediates between the compensated photon/null fix-sector and the
effective mass readout of closure-stable localized branches.
\end{theorem}

\begin{proof}
By the previous null-axis analysis, the photon sector is the terminal
compensated fix-sector. By the mass analysis, massive readout is residual
localization relative to that fix-sector. Hence the remaining structural task
is to identify the admissible interface through which this residual
localization is expressed in the physical readout layer. The Higgs sector is
precisely this interface.
\end{proof}

